% A4-stav-dp.tex -- dvojstranový podklad na konzultáciu.
%
% Čísla sa neprepisujú: načítavajú sa z tých istých generovaných makier, z
% ktorých ich berie práca. Preklad tohto súboru teda nemôže ukazovať inú
% hodnotu než PDF práce.
%
%   cd ai-doc && pdflatex -interaction=nonstopmode A4-stav-dp.tex
\documentclass[9.5pt,a4paper]{article}

\usepackage[utf8]{inputenc}
\usepackage[T1]{fontenc}
\usepackage[slovak]{babel}
\usepackage{lmodern}
\usepackage[margin=16mm,top=14mm,bottom=14mm]{geometry}
\usepackage{microtype}
\usepackage{enumitem}
\usepackage{booktabs}
\usepackage{tabularx}
\usepackage{xcolor}
\usepackage{titlesec}
\usepackage[hidelinks]{hyperref}

\definecolor{ink}{HTML}{1C1F26}
\definecolor{accent}{HTML}{2B4C7E}
\definecolor{warn}{HTML}{8A4A1C}
\definecolor{rule}{HTML}{C9C4B8}
\color{ink}

\titleformat{\section}{\normalfont\bfseries\large\color{accent}}{}{0pt}{}
\titlespacing{\section}{0pt}{10pt}{4pt}
\titleformat{\subsection}{\normalfont\bfseries}{}{0pt}{}
\titlespacing{\subsection}{0pt}{7pt}{2pt}

\setlist[itemize]{nosep,leftmargin=12pt,topsep=2pt}
\setlist[enumerate]{nosep,leftmargin=16pt,topsep=2pt}
\setlength{\parindent}{0pt}
\setlength{\parskip}{3pt}
\renewcommand{\arraystretch}{1.1}

% Generated numbers, from the same two files the thesis loads.
\input{../mastersLaTeX/analysis/out/thesis_numbers}
% Generated by assessment/tex_export.py -- do not edit.
% Numbers about the measuring instruments themselves.

\newcommand{\instrItemsTotal}{96}
\newcommand{\instrBanks}{8}
\newcommand{\instrFormItems}{12}
\newcommand{\instrFormPoints}{15}
\newcommand{\instrAnchors}{6}
\newcommand{\instrOpenPerForm}{3}
\newcommand{\instrCtItems}{12}
\newcommand{\instrMaiItems}{18}
\newcommand{\instrPrecheckItems}{5}
\newcommand{\instrTtItems}{6}
\newcommand{\instrTtPoints}{18}
\newcommand{\instrMeasureLessons}{4}
\newcommand{\instrBudgetUnderstated}{25.6}
\newcommand{\instrBudgetWorst}{TT}
\newcommand{\instrBudgetWorstBy}{18.5}
\newcommand{\instrLessonMin}{64}
\newcommand{\instrFatigueCap}{45}
\newcommand{\instrTests}{117}
\newcommand{\instrTestSections}{14}


\pagestyle{empty}

\begin{document}

{\LARGE\bfseries Ako nástroj overuje a vyhodnocuje žiakov}\\[2pt]
{\large\color{warn} Čo pribudlo od poslednej konzultácie a čo sa pri tom našlo}\\[4pt]
{\small\color{warn} Stav k 9.~októbru 2026 \quad\textbullet\quad 221 strán \quad\textbullet\quad
\repoTests{} testov v \repoSuites{} sadách \quad\textbullet\quad \repoChecks{} konzistenčných
kontrol \quad\textbullet\quad \repoMacros{} generovaných hodnôt \quad\textbullet\quad
\textbf{zber v školách neprebehol}}

\vspace{2pt}\textcolor{rule}{\hrule}\vspace{4pt}

\section{Medzera, ktorá sa zatvorila}

Nástroj vedel, \emph{ako sa žiak pýta}: záchytná vrstva zbierala prompty,
klasifikátor im priradil úroveň, index PEQI z nich spočítal epistemickú kvalitu
interakcie. O žiakovi samom nevedel nič~--- testy existovali len ako text
prílohy, nikto ich nepodával, neskóroval ani nekontroloval. Zadanie pritom žiada
\emph{meranie zmien v schopnostiach pred a po použití nástroja}, a zmena je
rozdiel dvoch meraní tej istej osoby. Pribudol hodnotiaci modul
(\texttt{assessment/}, štandardná knižnica, \instrTests{} testov v
\instrTestSections{} oddieloch), jeho overenie na simulovaných dátach
(\texttt{analysis/assessment\_validation.py}), nová časť metodiky a príloha~H.

\subsection{Zásada: položka je údaj, nie próza}

Zásobník \texttt{assessment/banks/*.json} (\asmBanks{}~nástrojov,
\asmItems{}~položiek, \asmDistractorsDiagnosed{}~diagnóz rozptyľovačov) je
jediný zdroj pravdy. Čítajú z neho tri programy: ten, čo položky žiakovi
podáva, ten, čo ich skóruje, a ten, čo tlačí prílohu~A. \textbf{Príloha sa od
nástroja rozísť nemôže, pretože z neho vzniká}~--- preklad práce zlyhá, ak sa
vytlačené zadania prestanú zhodovať so zásobníkom. Dve kópie jednej
špecifikácie sa v tomto projekte rozišli už dvakrát a oprava vždy šla len do
jednej.

\section{Čo sa pri tom našlo}

\subsection{1. Kotviace položky neboli totožné. Chyba obracia znak.}

Metodika aj príloha tvrdili, že \emph{šesť položiek je v oboch testových formách
totožných}, a že práve tieto kotviace položky robia formy porovnateľnými.
Totožná nebola ani jedna: okamžitá forma sa pýtala na pole s 1024~prvkami,
odložená na pole s 512; jedna „kotviaca“ položka mala v jednej forme štyri
možnosti a v druhej voľnú numerickú odpoveď. Kotviaca položka s iným $n$ je
\emph{iná} položka s inou obtiažnosťou, takže prevod oboch foriem na spoločnú
škálu neexistoval.

\begin{center}\small
\begin{tabularx}{\textwidth}{@{}lccX@{}}
\toprule
\textbf{Rozdiel kotviacich položiek} & \textbf{Odhad} & \textbf{Skreslenie} &
\textbf{Čo by práca vykázala} \\
\midrule
žiadny (opravený stav) & \asmEqEstZero  & \asmEqBiasZero  & zabúdanie, správne \\
0{,}3 logitu           & \asmEqEstThree & \asmEqBiasThree & zabúdanie, výrazne podhodnotené \\
\textcolor{warn}{\textbf{0{,}5 logitu}} & \textcolor{warn}{\textbf{\asmEqEstHalf}} &
\textcolor{warn}{\textbf{\asmEqBiasHalf}} &
\textcolor{warn}{\textbf{zlepšenie namiesto zabúdania}} \\
0{,}8 logitu           & \asmEqEstEight & \asmEqBiasEight & zlepšenie, veľké \\
\bottomrule
\end{tabularx}
\end{center}

Do simulácie sa vložilo zabúdanie \asmEqPlanted{}~logitu. Pri pol logitu
rozdielu~--- čo je pri 1024 oproti 512 striedmy odhad~--- by štúdia vykázala, že
žiaci sa \emph{zlepšili}. A nič by na to neupozornilo: model konverguje, zhoda
položiek s modelom prejde, interval je úzky. \textbf{Chyba, ktorá nevyzerá ako
chyba}, je v metodologickej práci najhorší možný druh.

\textbf{Oprava:} obe formy majú \asmFormItems{}~položiek, jednu špecifikačnú
tabuľku, \asmOpenPerForm{}~rubrikové položky a \asmAnchors{} znak po znaku
totožných kotviacich položiek, výhradne uzavretých. Validátor zásobníka to
vynucuje. Cena je spoľahlivosť $0{,}73$ namiesto $0{,}78$~--- rovnako ako
odložená forma, ktorá je primárnym výstupom.

\subsection{2. Oprava má vlastnú cenu a tá je väčšia, než sa dalo tušiť}

Ak sú kotviace položky totožné, žiak ich v odloženom teste vidí druhýkrát, a to
je vybavovacie precvičenie. Pri precvičení $0{,}25$~logitu sa vložené zabúdanie
\asmEqPlanted{} odhadne na \asmTeForgetting{}~--- takmer celé zmizne. Kontrast
medzi podmienkami je však chránený: \asmTeContrast{} pri vloženom
\asmTeContrastPlanted{}. \textbf{Konfirmačná rodina je v poriadku, absolútna
úroveň zadržania nie}, a to sa vykazuje ako priznané obmedzenie.

\textbf{Rotované kotvenie} to vie opraviť: pre každého žiaka sa jeho šesť
kotviacich položiek vyberie \emph{náhodne} zo zásobníka \asmRotPool{}~položiek
na tému. V odloženom teste potom na každú položku odpovedá aj žiak, ktorý ju
videl, aj žiak, ktorý nie, a tie dve skupiny sú pridelením ekvivalentné.
Precvičenie sa odhadne na \asmRotPractice{} (vložené \asmRotPracticePlanted{}),
zabúdanie na \asmRotForgetting{} oproti \asmFixForgetting{} pri pevnej sade.
Rotácia zmenšuje skreslenie približne trojnásobne. \textbf{Je to odporúčaná
zmena s pilotom, nie zmena vykonaná}: zásobník treba najprv nakalibrovať.

\subsection{3. Časový rozpočet: tri z piatich stretnutí sa nezmestili}

Rozpočet sa počítal ako počet položiek krát \emph{priemerný} čas na položku. To
je v poriadku, keď položky stoja rovnako, a tu nie: okamžitý test mal desať
uzavretých položiek po 40~sekundách a šesť rubrikových po dvoch a pol minúte;
transferovému testu, ktorému príloha sama priznáva 25~minút, dal priemer šesť a
pol. Chyby sa nevyrušia, pretože ukazujú rovnakým smerom: každý podhodnotený
nástroj je nástroj s rubrikovými položkami, a tie sedia v stretnutí~T4, ktoré
nesie primárny výstup. Podhodnotenie \asmBudgetUnderstated{}~minúty celkom,
najviac transferový test o \asmBudgetWorstBy{}.

Štyri zmeny: okamžitá forma \asmFormItems{}~položiek; predtest k téme na konec
\emph{predchádzajúcej} hodiny; výklad \instrLessonMin{}~minút namiesto 75
(kontrolný list vernosti sčítal 90~minút vyučovania v 90-minútovom bloku, ktorý
má obsahovať aj test~--- rozvrh a kontrolný list si protirečili o pätnásť minút);
a \textbf{T4 sa delí na T4a a T4b}. Pribudol únavový strop
\instrFatigueCap{}~minút na najdlhší \emph{neprerušený} úsek: únava sa násobí s
motiváciou, ktorú podmienky ovplyvňujú, takže priveľmi dlhé stretnutie nepridáva
šum, ale šum korelovaný so zásahom. \textbf{Škola dá \asmMeasureLessons{}
meracie hodiny namiesto troch.}

\subsection{4. Preklep v kóde žiaka nevyrobí chybu, ale dve}

Kód sa pri 330~žiakoch a piatich príležitostiach píše asi 1\,650-krát. Preklep
spôsobí, že jeden žiak ticho stratí príležitosť~--- a s ňou skóre zmeny, teda
veličinu, pre ktorú štúdia existuje~--- a žiak, ktorý neexistuje, ticho získa
riadok, ktorý analytický plán doplní imputáciou. Tri vrstvy: kontrolný znak
(ISO~7064), abeceda bez zámeniteľných znakov (nástroj povie, \emph{ktorá} chyba
sa stala), a „jeden kód, jedno sedenie“~--- jediná vrstva, ktorá zachytí žiaka
zadávajúceho \emph{platný} kód niekoho iného.

Norma tvrdí, že zachytí každé prehodenie susedných znakov; pri vyčerpávajúcom
preverení prejde jedno z dvoch tisíc. Kódový priestor má $9\,765\,625$~prvkov a
škola potrebuje tridsať, takže sa vydávajú len kódy, pri ktorých to neplatí:
\textbf{100\,\% konštrukciou, nie 99{,}95\,\% citáciou.}

\subsection{5. Dve veci, ktoré boli doteraz tvrdením}

\textbf{Úsilie pri teste.} „Žiaci sa naučili viac“ a „žiaci sa viac snažili na
teste po hodine, ktorá ich bavila“ sú ten istý kontrast, pokiaľ ich niečo
neoddelí. Podmienka~$G$ odstraňuje úsilie počas hodiny, takže rozdiel v úsilí je
pravdepodobne diferenciálny podľa podmienky, teda skreslenie a nie šum. Rozdiel
$0{,}08$, taký, ktorý by vedel predstierať najmenší zaujímavý účinok, sa odhalí
v \asmEffortDetect{} prípadov pri falošnej pozitivite \asmEffortFalsePos{}.
Papierový test toto číslo nevyrobí vôbec.

\textbf{Kontaminácia testu.} Doterajšou odpoveďou bola položka V12 kontrolného
listu~--- tvrdenie dozorcu o dvadsiatich obrazovkách. Nástroj zaznamenáva päť
rodín príznakov z troch nezávislých zdrojov, vylúči odpoveď len pri dvoch
príznakoch z \emph{rôznych} rodín, a je kalibrovaný na stretnutí~T0a, kde je
kontaminácia nemožná: čokoľvek tam označí, je falošná pozitivita konštrukciou.
Pri realistickom nazeraní (dve až štyri položky z dvanástich) zachytí
\asmIntegDetect{} kontaminovaných žiakov pri falošnej pozitivite
\asmIntegFalsePos{}. Žiadny žiak sa škole nepomenuje: príznak, ktorý by sa mohol
dotknúť známky, by zmenil motiváciu písať test poctivo a zničil by meranie, pre
ktoré bol zavedený.

\subsection{6. Kategorizácia, ktorú vie použiť aj jedna trieda}

Latentná profilová analýza je tvrdenie o \emph{vzorke}: potrebuje stovky žiakov
a označenie jedného žiaka závisí od toho, kto je v súbore s ním. Učiteľ s jednou
triedou z nej nedostane nič~--- takže zadanie o kategorizácii používateľov bolo
splnené spôsobom, ktorý vie použiť len štúdia. Nástroj preto obsahuje aj
deterministické pravidlo nad \emph{absolútnymi}, vopred stanovenými hranicami,
odvodenými z významu škály (PEQI $0{,}25$ a $0{,}50$ sú hranice $L_1$ a $L_2$
samotnej taxonómie).

\begin{center}\small
\begin{tabular}{@{}lccc@{}}
\toprule
\textbf{Počet žiakov} & \textbf{Pravidlo} & \textbf{Zmesový model} & \textbf{Komponentov} \\
\midrule
\textbf{\asmNClass{} (jedna trieda)} & \textbf{\asmRuleAriClass} &
\textcolor{warn}{\textbf{\asmLpaAriClass}} & \asmLpaKClass \\
\asmNThree{}               & \asmRuleAriThree & \asmLpaAriThree & \asmLpaKThree \\
\asmNStudy{} (táto štúdia) & \asmRuleAriStudy & \textbf{\asmLpaAriStudy} & \asmLpaKStudy \\
\asmNLarge{}               & \asmRuleAriLarge & \asmLpaAriLarge & \asmLpaKLarge \\
\bottomrule
\end{tabular}
\end{center}

Pri jednej triede zmesový model nenájde žiadne profily; pri \asmNStudy{} a viac
je lepší než pravidlo. \textbf{Nie sú to konkurenti}: zmes je výskumná otázka o
populácii, pravidlo je to, čo dostane škola. Ten istý žiak preklasifikovaný v
\asmLabelCohorts{}~kohortách dostal vždy to isté označenie (\asmLabelStable).
Posun jednej hranice o $\pm$20\,\% mení zhodu o \asmCutSingle, všetkých piatich
naraz o \asmCutJoint{}~--- pravidlo \emph{je} citlivé na ich spoločné nastavenie,
a preto musia byť predregistrované. Test proti svetu, v ktorom je profil len
preznačkovaná vstupná úroveň, má silu \asmIncPower{} pri falošnej pozitivite
\asmIncFalsePos{}.

\subsection{7. Čo sa vracia škole}

Učiteľ dostane počty mylných predstáv po položkách~--- nie „jedenásť žiakov nemá
bod“, ale „jedenásť žiakov si myslí, že zdvojnásobenie poľa zdvojnásobí počet
krokov“~--- bez rozpisu podľa podmienky, ktorý mu neslúži a vyrába bunky s dvoma
žiakmi. Škola dostane agregáty s potlačením buniek pod piatimi žiakmi. Žiak
nedostáva individuálne skóre: spoľahlivosť nástrojov unesie porovnanie skupín,
nie výrok o jednotlivcovi.

\section{Čo potrebujem rozhodnúť}

\begin{enumerate}
    \item \textbf{Rotované kotvenie:} pilotovať, alebo priznať obmedzenie?
    Odporúčam priznať teraz, pilotovať rotáciu, rozhodnúť pred zberom.
    \item \textbf{Štvrtá meracia hodina:} unesie to nábor? Alternatívou je
    skrátiť transferový test, teda práve to meranie, na ktorom má byť rozdiel
    podľa hypotézy najsilnejší.
    \item \textbf{Predregistrácia piatich hraníc pravidla} do štatistického plánu
    (príloha~C), skôr než existujú akékoľvek dáta.
    \item Nezmenené: školy \textbullet{} etická komisia (kritická cesta)
    \textbullet{} pilot alebo plná štúdia \textbullet{} autorstvo citácie
    Hariyanto (2025) \textbullet{} má zmysel udržať fyzikálne rameno?
\end{enumerate}

% The footer is two lines and belongs on the last text page, not alone on a new
% one. Three lines of extra height on this page is cheaper than cutting content.
\enlargethispage{3\baselineskip}
\vspace{2pt}\textcolor{rule}{\hrule}\vspace{3pt}
{\footnotesize\color{warn} Žiadne číslo v tomto dokumente nie je napísané rukou.
Všetky sa načítavajú z tých istých generovaných makier, z ktorých ich berie
práca, takže tento podklad nemôže ukazovať inú hodnotu než jej PDF.
Rozhodnutia a ich dôvody sú v \texttt{decisions.md}.}

\end{document}
